\section{Introduction}\label{sec:introduction}

When a civil registry, a land registry, and an employer each issue digital credentials that can be used for a government housing subsidy, covering family status, property records, and income respectively, the bank processing the application must accept all three in a single eligibility check. EU regulation mandates a specific credential format for government attestations \citep{eidas2}. Data protection law requires that privacy-sensitive checks disclose only the minimum necessary information \citep{gdpr}; an income threshold check should not require revealing the exact value, yet the mandated format lacks this capability while other standardized formats provide it. No existing modeling tool is available to detect this and similar requirement conflicts. These credentials are tamper-evident, cryptographically verifiable claims exchanged between issuers, holders, and verifiers as defined by the W3C Verifiable Credentials Data Model 2.0 \citep{sporny_verifiable_2025}; designing their ecosystems requires satisfying constraints that span what claims mean, how credentials are structured, and what each format can express.

Credential ecosystem designs can contain errors undetectable by single-layer inspection because governance frameworks impose requirements that conflict across abstraction layers \citep{mazzocca_survey_2025}. The W3C Verifiable Credentials Data Model defines structural conformance rules that standard credentials must satisfy \citep{sporny_verifiable_2025}. The eIDAS Architecture and Reference Framework mandates specific credential formats for government-issued attestations \citep{noauthor_eu-digital-identity-walleteudi-doc-architecture-and-reference-framework_2026}. The General Data Protection Regulation requires data minimization \citep{gdpr}, operationalized in this paper as a format capability requirement: a threshold comparison should not require disclosing the underlying value. Unlike hierarchical requirement systems where constraints decompose top-down, credential governance sources are non-cooperating peers with competing goals; their requirements cannot necessarily be brought into a coherent whole.

Existing tools operate at a single layer: JSON Schema validators check credential structure, format-specific conformance tools verify encoding constraints, governance frameworks define requirements in isolation. None checks cross-layer consistency. In the housing subsidy scenario, the income credential's format assignment, structural conformance, and privacy requirement each pass their respective single-layer checks, yet their combination is unsatisfiable. A credential schema may be well-formed when inspected in isolation, yet violate a cross-layer constraint that links domain-level claim semantics to format-specific privacy capabilities. Detecting such errors requires a formalization that spans all three concern spaces.

We make the following contributions:

\begin{enumerate}
\def\labelenumi{\arabic{enumi}.}
\tightlist
\item
  A \textbf{three-layer metamodel} for credential ecosystem design, grounded in the W3C Verifiable Credentials Data Model 2.0, spanning claim properties, credential schemas, and format-specific representations.
\item
  A \textbf{formalization of cross-layer constraints} as predicates in the Refinery partial graph modeling framework.
\item
  A \textbf{three-axis validation} demonstrating metamodel coverage against the W3C specification, constraint expressiveness against W3C and EU regulatory sources, and error detection against known credential design anti-patterns. The validation detects two cross-layer design errors, a governance conflict and a format expressiveness gap, undetectable by single-layer analysis.
\end{enumerate}

Cross-layer constraints are formalized as graph predicates in Refinery \citep{marussy_refinery_2024}. Refinery's solver reveals when constraints from different governance frameworks are formally contradictory and generates valid alternative designs when they are not. The framework supports three usage modes: \emph{consistency checking} confirms that a design satisfies all constraints, \emph{error identification} names specific violations, and \emph{design space exploration} generates valid configurations or proves that none exist. \autoref{fig:teaser} illustrates the complete pipeline on a housing subsidy credential scenario that serves as the running example throughout.

\autoref{sec:background} introduces the W3C \ac{VCDM} 2.0, multi-layer modeling, and the Refinery graph modeling tool, which provides the formal apparatus for the cross-layer analysis. \autoref{sec:overview} develops a housing subsidy scenario in which credentials from three independent authorities fall under conflicting governance mandates, and defines the three usage modalities that set the requirements the formalization must satisfy. \autoref{sec:approach} formalizes the three-layer metamodel and its cross-layer constraints as Refinery graph predicates grounded in the running example. \autoref{sec:evaluation} validates the metamodel against the W3C specification, EU regulatory sources, and solver scalability. \autoref{sec:related-work} positions the contribution against credential ecosystem formalization and model-driven security engineering, and \autoref{sec:conclusion} examines the design-time scope boundary and identifies extensions to runtime verification.
